%\documentclass[12pt]{scrreprt}
\documentclass[12pt]{report} 

% language may be romanian or english (default is english)
% type may be bachelor or master (default is bachelor)
\usepackage[language=english, type=bachelor]{style}

%\geometry{a4paper,top=2.5cm,left=3cm,right=2.5cm,bottom=2.5cm}
%in style
%controlling the appearance of your headers and footers
\usepackage{fancyhdr}
\pagestyle{fancy}
\lhead{}
\chead{}
\renewcommand{\headrulewidth}{0.2pt}
\renewcommand{\footrulewidth}{0.2pt}

\begin{document}

\specialization{Mathematics and Computer Science}	
\title{Agent-Based Modelling of Brown Bear Population Dynamics in the Romanian Carpathians}					   
\author{[Student name]}											
\supervisor{[Degree, title and supervisor name]}	
				
\maketitle


\newpage
\thispagestyle{empty}
\mbox{}
\newpage
\pagenumbering{roman} 

\cleardoublepage
ABSTRACT
\vspace{0.5cm}	
\hrule
\vspace{0.5cm}	
%\cleardoublepage

Brown bear population dynamics in Romania represent both an ecological and a public-policy challenge. This thesis develops a spatial agent-based simulation focused on a Transylvania case study inside the Romanian Carpathians and uses it to investigate the population increase observed in the 2005--2021 interval. The objective is explanatory and counterfactual: to compare plausible driver combinations under explicit assumptions, not to produce exact deterministic forecasts.

The model represents each bear as an autonomous agent with age, sex, satiety, reproduction, and movement behaviour, while the environment is represented as a heterogeneous grid with food and danger dynamics. The simulation architecture separates a parallel planning phase from a deterministic commit phase, and is supported by a reproducible experiment pipeline including seeded runs, schedule-driven parameter updates, and per-tick metrics recording.

The thesis combines ecological validation and computational validation. Ecological results are evaluated through baseline and counterfactual scenarios, using trajectory-fit and population-structure indicators. In the final 2005--2021 campaign, the best fit is obtained by the anthropogenic-food schedule (MAE 261.25, RMSE 336.27), while baseline and policy-only schedules remain substantially less accurate. Computational results compare sequential, multithreaded, and distributed execution modes. A central technical finding is that for this fine-grained tick-synchronized model, distributed execution inside a single run can be dominated by communication and synchronization overhead, while distributed execution is effective for large batches of independent runs.

Overall, the contribution is twofold: a practical ABM framework for hypothesis testing in brown bear population dynamics, and an evidence-based execution strategy that aligns model granularity with computational design.



\tableofcontents


\newpage
\pagenumbering{arabic}

\chapter{Introduction}
\label{ch:introduction}

\section{Context and motivation}

In the last years, the brown bear population in Romania has become one of the most visible and debated wildlife-management topics in Europe. The debate is active in public media, policy documents, and scientific reports. On one side, there are concerns related to safety and conflict at settlement boundaries. On the other side, there is the scientific need to explain long-term population dynamics using measurable mechanisms instead of isolated observations.

The difficulty is that many mechanisms can act at the same time: hunting-policy changes, supplemental feeding, anthropogenic food near villages, landscape transition, and seasonality. If these mechanisms are analyzed separately, interpretation may become incomplete. In this context, simulation is useful because it allows controlled experiments where each factor can be switched, scheduled, and measured in repeated runs.

This project proposes a spatial agent-based simulation for brown bears, inspired by the Romanian Carpathian context and implemented in Java. The approach aims to preserve biological meaning while remaining computationally practical for large experiment batches.

\section{Research problem and objectives}

The core research question is the following: \emph{which mechanisms, alone or in combination, can best explain the observed increase of brown bear population in the Romanian Carpathian area during 2005--2021?}

The objectives of this thesis are:

\begin{enumerate}
    \item to design an agent-based model where each bear is represented as an individual with age, sex, satiety, reproduction and movement behaviour;
    \item to represent habitat and risk as spatially heterogeneous grid cells with food and danger dynamics;
    \item to evaluate candidate drivers using baseline and counterfactual scenarios;
    \item to build a reproducible experiment pipeline based on controlled random seeds, metrics recording, and parameter schedules;
    \item to evaluate computational behaviour with emphasis on multithreaded versus distributed execution strategy.
\end{enumerate}

\section{Scope and delimitations}

This thesis adopts a focused and explicit scope. The study is defined as a \textbf{Transylvania-focused case study inside the Romanian Carpathians}, with a \textbf{2005--2021 analysis window}.

The purpose of the model is \textbf{explanatory counterfactual analysis}, not exact numerical forecasting. In other words, the simulation is used to compare how plausible driver combinations can reproduce the observed increase trend under explicit assumptions, rather than to claim exact future population values.

This scope is chosen for three reasons:

\begin{enumerate}
    \item it matches the model granularity (local interactions, tick-based dynamics, scenario schedules);
    \item it is scientifically more robust under real-world uncertainty in census and policy data;
    \item it aligns with the technical contribution, where many controlled scenario replicates are required.
\end{enumerate}

Therefore, conclusions in this thesis should be interpreted as comparative explanatory evidence across scenarios, not as deterministic prediction statements.

\section{Original contributions}

The main contributions are both ecological and technical, and they are tightly connected.

First, the ecological contribution is the integration of multiple candidate drivers inside one simulation framework. Instead of discussing one factor at a time, the project evaluates both isolated and combined effects. This allows interpretation of relative contribution and interaction effects.

Second, the technical contribution is an execution strategy that is realistic for this model class. The implementation supports multithreading and distributed execution, but the thesis does not assume distributed is always faster. For this fine-grained tick-based simulation, distributed execution \emph{inside one run} can introduce communication and synchronization overhead larger than the useful work. In contrast, distributed execution is expected to be efficient for \emph{batches of independent runs}.

Third, the project includes practical reproducibility components directly in code: run configurations, seeded randomness, per-tick metrics, schedule-based parameter updates, and calibration utilities.

Fourth, the thesis contributes a clear modelling scope statement that keeps claims proportional to evidence: explanatory counterfactual analysis over 2005--2021 for a Transylvania-focused Carpathian case study.

\section{Thesis structure}

The thesis is organized in five main chapters plus references.

Chapter 2 presents ecological background and related work in ABM and parallel simulation, together with the specific gap addressed by this project. Chapter 3 describes the proposed approach, including model structure, agent logic, environment representation, and execution architecture. Chapter 4 presents the evaluation design and templates used for reporting ecological and computational results. Chapter 5 discusses findings, limitations, conclusions, and future work.

To reach the target manuscript size of approximately 30 pages in the current template format, the intended distribution is:

\begin{itemize}
    \item Chapter 1: 4--5 pages,
    \item Chapter 2: 7--8 pages,
    \item Chapter 3: 8--9 pages,
    \item Chapter 4: 6--7 pages,
    \item Chapter 5: 3--4 pages.
\end{itemize}

This page budget is used as a writing guide and will be refined after figures, tables, and final numerical results are inserted.

%\addcontentsline{toc}{chapter}{Introducere}
%\addcontentsline{toc}{chapter}{Introduction}

\chapter{Background and related work}
\label{ch:background_related}

\section{Brown bear ecology and demographic context}

Brown bear ecology is a very good case for simulation, because the population trends we read about in reports are never the product of a single cause. They emerge from many local processes acting at the same time: reproduction, mortality at different ages, access to food, the way individuals use space, and the constant friction of living next to humans. In practice none of these processes acts alone. They overlap in time and they interact, and their combined effect is what pushes a population toward growth, decline, or a local equilibrium. This is precisely why a verbal, one-cause-at-a-time explanation tends to fall short, and why a model that lets the causes act together is useful.

From the demographic literature we know that brown bear trajectories are shaped above all by age and sex structure, by the long intervals between successive litters, by cub survival, and by the mortality pressure that human-dominated landscapes add on top of natural mortality \cite{Schwartz2003,Saether1998,Swenson1997}. These are not abstract details for a modeller; they are the realistic bounds that keep the model honest. When the simulation in this thesis uses a minimum reproductive age of four years, a gestation of about a year, a multi-year cooldown between litters, a litter mean close to two cubs, and elevated mortality for cubs, none of those numbers is invented. They come from this body of work, and the role of the literature is exactly to provide defensible ranges instead of arbitrary values \cite{Swenson2000}.

The Romanian context makes the problem more relevant still. Romania hosts by a wide margin the largest brown bear population in the European Union, and recent work has documented both the movement ecology of these bears and the patterns of human--bear interaction in the Carpathians \cite{Pop2018,Popescu2017}. At the same time, institutional reports and policy documents point to substantial changes both in the estimated size of the population and in the management strategy applied to it over the last two decades \cite{INCDS2023,MMAP2024}. The increase in conflict reports and the shifting estimates of population size are not just background colour; they are the very signal this thesis tries to explain. So this is not only a theoretical modelling exercise. It is a problem with direct public-policy weight, where the question of which mechanism actually drove the rise has practical consequences for how the population should be managed \cite{Chapron2014}.

\section{Simulation and the agent-based paradigm}

Agent-based modelling is the natural tool to reach for when individual heterogeneity matters and when interactions are local in space and time, and both of those are true for bears. Instead of forcing the whole population into one aggregated equation, an agent-based model lets each individual carry its own internal state and update its behaviour according to its own local context. The population-level behaviour is then not assumed in advance; it appears as an emergent consequence of all those individual updates. For a question like ours, where we want to test whether a given assumption is \emph{sufficient} to produce an observed trend, this emergence is the whole value of the method.

The methodological foundations of individual-based ecological models are by now well established \cite{GrimmRailsback2005,RailsbackGrimm2019,DeAngelis2005}. Just as importantly, the field has developed shared standards for describing such models, most notably the ODD protocol and its update, which lay out a practical structure for stating a model's entities, processes, scheduling and reproducibility decisions \cite{Grimm2006ODD,Grimm2020ODD}. This thesis follows the same spirit, even where it does not follow the protocol to the letter: Chapter~3 explicitly defines the agent state, the scheduling of the tick loop, the environment representation and the experiment protocol, precisely so that the model can be understood and, in principle, reproduced.

The reason this paradigm fits counterfactual analysis so well is that interventions can be expressed as explicit rules and then switched on and off. If we want to test the effect of changing hunting pressure or of adding a food subsidy, we do not have to re-derive an equation. We change a rule or a parameter at a given year and re-run, and because the model is stochastic, we re-run it several times and look at the distribution of outcomes rather than a single line. This is the same logic used across the wider agent-based modelling community, from general treatments of the method \cite{Macal2016} to teaching-oriented platforms such as NetLogo \cite{WilenskyRand2015}.

\section{Romanian case study framing}

In this thesis the Romanian situation is treated as a concrete testbed, not as a generic backdrop. More precisely, the empirical framing is the brown bear population across Romania's Carpathian forests as a whole, with 2005--2021 as the main analysis window. Choosing a real framing rather than a purely abstract one brings two clear advantages.

First, it provides real-world constraints. Demographic assumptions, habitat structure and the policy timeline can all be tied to published or institutional sources, which sharply reduces the number of arbitrary modelling choices and makes the scenarios defensible. The landscape transitions of the Romanian Carpathians over recent decades have themselves been studied \cite{Munteanu2014,Griffiths2013}, which gives some grounding to the habitat-loss scenario we test later. Second, a real framing allows the results to be interpreted in a management-oriented direction. Even though the model stays simplified, its scenario outputs can still be read as statements like ``this combination of pressures reproduces the observed trend more plausibly than that one'', which is the kind of statement a manager can actually use.

At the same time, this framing brings a responsibility in how the results are read. Institutional population estimates evolve as census methods improve, and policy periods do not always begin and end on clean calendar boundaries. For that reason the simulation is presented as a structured comparative tool under explicit assumptions, and not as an exact forecasting device. I would rather make a careful comparative claim that holds than a precise numerical claim that the data cannot support.

\section{Demography and the problem of parameter realism}

There is a recurring weakness in student simulation projects that I was determined to avoid, and it is worth naming directly: parameter arbitrariness. In ecological agent-based modelling it is dangerously easy to produce plots that look plausible but that mean very little, because the parameters behind them were tuned by hand until the picture looked nice. A model like that can be made to show almost anything, which means it shows nothing.

The defence against this is to constrain parameter ranges with the demographic literature wherever possible, and only then to calibrate within those ranges. The references used here support realistic intervals for the quantities that matter most: age-related mortality, the vulnerability and survival bottleneck of cubs, the intervals between litters, and the additional mortality that human pressure imposes \cite{Schwartz2003,Saether1998,McLellan1989,Bischof2009}. The sex- and age-specific mortality multipliers in the model, the elevated cub danger, and the long reproductive cooldown all sit inside ranges drawn from this work. This does not make the uncertainty disappear, but it does give a defensible starting point for the calibration and sensitivity experiments, which is a very different situation from picking numbers until the curve fits.

\section{Methodological references for evaluation}

Besides the ecological literature, the thesis also leans on methodological sources, because how you evaluate a stochastic simulation is itself a methodological question. Sensitivity analysis and parameter-estimation workflows for this kind of model are well covered in the standard references \cite{Saltelli2008,Thiele2014}, and they inform the way Chapter~4 reports its results. The decision that follows from them is simple but non-negotiable: if a claim depends on a stochastic simulation, a single run is not evidence. Every scenario in this thesis is therefore run as several seeded replicates and then aggregated, and the comparison against the reference trajectory is made on those aggregates rather than on one lucky or unlucky line.

\section{Related software approaches}

From the software side there is a mature body of prior work on multi-agent simulation environments, including MASON and Repast and their high-performance, distributed variants \cite{Luke2005,CollierNorth2013,North2013}, lighter frameworks such as Mesa \cite{Masad2015}, and dedicated spatially explicit ecological platforms \cite{Bocedi2014}. There is also a long-standing theoretical foundation for parallel and distributed simulation in general \cite{Fujimoto2000}. I chose to build the model directly in Java rather than on top of one of these frameworks, partly for full control over the execution and reproducibility machinery, and partly because the computational question I wanted to ask is about that very machinery.

What this literature makes clear, and what is central to my argument, is that the benefit of parallel or distributed execution depends entirely on the shape of the workload: on how heavy the per-agent computation is, on the communication pattern, and on how often the model has to synchronise. In models where each agent does heavy local computation between synchronisations, distribution scales well. In models where the per-agent work is light and synchronisation is frequent, communication overhead can erode the speedup and even reverse it. This distinction is exactly the one that applies to the bear model. Because it is tick-based and globally synchronised, and because each agent action is cheap, distributing a single run can be slower than running it multithreaded on one machine, while distributing many independent runs remains very effective. The wider ecological agent-based literature, including work on other large carnivores such as tigers \cite{Carter2015}, confirms more generally that complex wildlife dynamics can be represented faithfully at the individual level and produce meaningful scenario outputs, which is the methodological reassurance this thesis needed before committing to the approach.

\section{Gap addressed by this thesis}

The gap this thesis addresses has two sides, one ecological and one computational, and the contribution is to close both at once with a single piece of work.

On the ecological side, much of the public and even technical discussion about bear-population change in Romania stays descriptive, and it tends to emphasise one driver at a time, it was the hunting ban, or it was the feeding, or it was habitat change, without ever testing how these drivers behave together inside one consistent framework. This thesis addresses that by integrating several candidate drivers into one model and evaluating them, both in isolation and combined, under controlled scenarios with replicated runs.

On the computational side, reports and student projects alike often present multithreaded and distributed execution as if they were universally beneficial, almost as a synonym for ``more serious''. In practice, as the simulation-systems literature warns, that is not true for every workload. For this tick-synchronised model the thesis works out explicitly where parallelism actually helps, and, just as importantly, where it would not. The model is therefore executed with a multithreaded, single-machine design, and a distributed execution is deliberately \emph{not} built, because the analysis shows it would not pay off for this granularity. The reasoned conclusion is a simple two-level picture:

\begin{itemize}
    \item multithreading is used to accelerate a single simulation run on one machine, where the threads share memory and the per-tick synchronisation barrier is cheap;
    \item distribution would only make sense at the level of large batches of independent runs, which are embarrassingly parallel, but at the scale of these experiments a single multi-core machine ran the whole campaign comfortably, so it was not needed and not implemented.
\end{itemize}

Stating this judgement explicitly, rather than reflexively reaching for a distributed implementation, improves both the scientific reproducibility of the results and the engineering honesty of how they were produced, and it is one of the things that I believe distinguishes this work from a model that simply runs and prints a curve.

\section{Chapter summary}

This chapter set out the conceptual and bibliographic foundation of the thesis. The ecological context motivates the modelling problem and supplies realistic parameter ranges; the agent-based literature supports the choice of paradigm and the standards for describing and evaluating such a model; and the simulation-systems literature supplies the criteria for the computational part of the evaluation. With this foundation in place, the next chapter moves from background to implementation: how the world is represented, how agents decide and act, how their effects are committed each tick, and how the whole thing is wrapped in a reproducible experiment pipeline.

\chapter{Proposed approach}
\label{ch:approach}

\section{Model overview}

The proposed solution is a spatial agent-based simulation written in Java, in which every brown bear in the studied area is represented as one individual software agent that lives, eats, moves, reproduces and eventually dies inside a shared, heterogeneous world. The choice of an agent-based model is not accidental. The phenomenon we are trying to understand, namely the increase of the Romanian brown bear population during 2005--2021, is the cumulative result of a very large number of small, local decisions: a single bear deciding whether to look for food, a single female deciding whether the conditions are good enough to reproduce, a single subadult male leaving the area where it was born. None of these decisions is interesting on its own, but together, repeated over many thousands of individuals and many simulated years, they produce the population trend that we actually observe in the field. The whole point of building the model this way is that the population trajectory is never written down anywhere in the code. It is an \emph{emergent} property, and this is exactly what we want, because it lets us ask whether a given combination of assumptions is able to reproduce the real trend without us forcing the answer by hand.

Time in the simulation is discrete and advances in equal ticks. One tick corresponds to a fixed fraction of a simulated year, and in the current calibration one tick is set to \texttt{0.0001140771} years, which is approximately one hour. This means a full simulated year is built out of roughly 8766 ticks, and the 2005--2021 window of sixteen years is therefore around 140{,}000 ticks long. Working at an hourly resolution may look excessive for a population-level question, but it is a deliberate decision: it lets behaviours such as feeding, daily movement and the slow decay of body condition play out at a believable pace, instead of being collapsed into one coarse yearly update where everything happens at once.

At each tick, the model follows a two-stage flow, which is one of the most important design decisions of the whole project:

\begin{enumerate}
    \item \textbf{Plan stage:} every agent looks at its own local surroundings, receives a percept, and decides on exactly one action for this tick. No agent is allowed to change the world during this stage. It only decides.
    \item \textbf{Commit stage:} all the decided actions are collected and then applied to the shared world in one controlled, ordered pass on a single thread.
\end{enumerate}

This separation between deciding and acting is what makes the rest of the architecture possible. Because nobody modifies the world while everybody is still thinking, the planning stage is free of data races and can be run in parallel across many threads without any locking on the world state. And because all the modifications are funnelled through one ordered commit, the resulting world transition is deterministic and easy to reason about. We will come back to this idea several times in this chapter, because almost every other design choice depends on it.

\section{Problem formulation and modelling assumptions}

Before describing the implementation in detail, it is worth stating clearly what kind of model this is and what it is not. The modelling objective is explanatory and counterfactual: we want to compare plausible explanations for the observed population increase under controlled, repeatable computational experiments. We are not trying to produce an exact deterministic forecast of how many bears Romania will have in a given future year, because the available census and policy data simply do not support that level of precision. What the model \emph{can} do well is answer a comparative question, of the form ``under these assumptions, does this driver, alone or combined with others, reproduce the observed trend better than that one?''.

In order to keep the model both interpretable and actually executable for large batches of runs, several assumptions are introduced, and it is worth listing them so the reader knows where the simplifications are:

\begin{itemize}
    \item the world is a regular two-dimensional grid, and interactions are local, restricted to a small neighbourhood around each agent;
    \item time advances in equal discrete ticks, with no continuous-time events between ticks;
    \item each agent perceives only its immediate surroundings, never the global state of the population;
    \item mortality and reproduction are probabilistic, but the probabilities are bounded by biological rules (minimum reproductive age, gestation, cooldown, litter limits);
    \item the environment is abstracted into food and danger fields layered on top of a small set of terrain categories.
\end{itemize}

None of these assumptions claims perfect realism, and I do not want to pretend otherwise. Their role is to strike a balance between biological plausibility on one side and reproducibility and computational tractability on the other. A model that is too detailed becomes impossible to calibrate and impossible to run thousands of times; a model that is too coarse stops being biologically meaningful. The assumptions above are the compromise that this thesis settled on, and the parameter values behind them are, wherever possible, tied to the demographic literature rather than chosen arbitrarily \cite{Swenson2000,Schwartz2003}.

\section{Agent state and behaviour}

Each bear agent is implemented in the \emph{BearAgent} class and carries a compact internal state. The fields that matter for behaviour are the bear's identity, its sex, its age in years, its satiety, its reproduction cooldown, its gestation progress, whether it is currently pregnant, and a remembered home-range anchor expressed as a pair of coordinates. It is important to notice that the agent does not store its own position on the map. The position lives in the world state, not in the agent, which keeps a clean separation between ``who the bear is'' and ``where the bear is'', and turns out to be very convenient when we save and restore whole populations later.

\subsection{The decision pipeline}

The heart of the agent is the method that decides the next action, and it is best described as a short pipeline that is evaluated in a fixed order at every tick. The order is meaningful, because the first condition that matches wins and the rest are skipped.

The very first thing that happens each tick is the passage of time. The agent ages by one tick, its reproduction cooldown and remaining gestation are both decremented, and its satiety decays. Only after time has passed do we check whether the bear has died. Death from starvation comes first: if satiety has reached zero, the agent decides to die with cause \texttt{STARVATION}. Next comes old age: once a bear is older than the configured maximum age, every tick it rolls against a small per-tick death probability, scaled by a sex-specific multiplier because in the model, as in reality, females are slightly longer-lived than males. After that comes danger: the danger value of the cell the bear currently stands on is turned into a per-tick probability of dying, and this probability is multiplied by several factors: a large multiplier for cubs, who are far more vulnerable, a sex multiplier because roaming males are more exposed, and a strong reduction during hibernation because a denned bear is sheltered.

If the bear survives all three mortality checks, the pipeline turns to reproduction and behaviour. A pregnant female whose gestation has finished gives birth, resets her cooldown and clears her pregnancy. If the bear is hibernating, the pipeline stops here: a denned bear does not forage, does not move and does not mate, although births can still happen in the den. Outside hibernation, an eligible female that perceives a male nearby becomes pregnant and starts her gestation clock. Then comes feeding: a hungry bear, meaning one whose satiety is below an eat-preference threshold, eats immediately if there is enough food on its current tile. If the bear is not hungry enough to eat in place, it evaluates its neighbouring cells and moves to the best one, provided that cell is actually better than staying put. And if there is no better neighbour but there is food where it stands, it eats anyway; otherwise it simply does nothing and waits for the next tick.

It is worth mentioning that this ordering is itself a modelling decision. By putting survival before reproduction and reproduction before feeding and movement, the agent expresses a simple priority: stay alive, then continue the species, then take care of daily needs. It is a deliberately reactive design rather than a planning one, because for a population-level question we care about the statistical behaviour of many bears, not about any single bear executing a clever long-term plan.

\subsection{Movement as a weighted trade-off}

Movement deserves its own discussion, because it is where most of the spatial realism of the model comes from. When a bear is not eating in place, it looks at the eight cells around it (a Moore neighbourhood) and gives each candidate cell a score. The cell with the highest score wins, and the bear moves there unless its current cell already scores best. The score is a weighted combination of four ingredients: the food on the cell pulls the bear in, the danger on the cell pushes it away, the local crowding (how many other bears are nearby) pushes it away, and a home-range penalty discourages the bear from wandering too far from the anchor point it remembers. A small amount of random noise is also added, so that bears do not behave like perfectly identical optimisers and the movement looks more natural.

The really interesting part is that the weights are not constant: they shift depending on how hungry the bear is. A well-fed bear weighs danger heavily and food lightly, so it behaves cautiously and stays away from risky cells such as villages and roads. A hungry bear does the opposite: the food weight grows, the danger weight shrinks, and the crowding and home-range weights are reduced as well. In plain words, a hungry bear is willing to take risks and to leave its comfortable range to find something to eat, which is exactly the kind of behaviour that pushes bears toward human settlements when natural food is scarce, and which is at the centre of the anthropogenic-food hypothesis we test in the next chapter. The hunger level is computed as a simple normalisation of satiety, and the weights are linearly interpolated between a ``when full'' value and a ``when hungry'' value, so the transition is smooth rather than a hard switch.

\subsection{Subadult male dispersal}

On top of the generic movement logic, the model adds a specific behaviour for subadult males, because in brown bears natal dispersal is overwhelmingly a male, subadult phenomenon \cite{McLellan1989,Steyaert2012}. A male bear inside a configured age window (roughly two to six years old) is treated as a disperser: its crowding weight is multiplied up so it actively avoids competition, its home-range weight is multiplied down so the anchor barely holds it, and it receives a positive bonus for every cell it places between itself and its birth range. The combined effect is that young males drift outward from where they were born, which spreads the population across the map instead of letting it pile up around the founding locations. This matters for the realism of the spatial pattern and, indirectly, for how the population finds and uses new food.

\subsection{Seasonality through hibernation}

Brown bears are strongly seasonal animals, and ignoring this would make the yearly dynamics wrong. The model represents seasonality with a hibernation window defined as a fraction of the year, configured to run from roughly early November to early March, wrapping around the year boundary. While a bear is inside this window, two things change. First, its metabolism slows down: satiety decays much more slowly, controlled by a hibernation metabolism multiplier well below one, so a denned bear can survive the winter on its reserves. Second, its exposure to danger drops sharply, because a sheltered, denning bear is far less likely to encounter the risks of the active season. During hibernation the agent does no foraging, no movement and no mating, but, importantly, births still occur, since brown bear cubs are in fact born in the den during winter. This single mechanism is responsible for a lot of the believable seasonal texture in the output.

\subsection{Reproduction and birth}

Reproduction in the model is gated by several biological constraints that all have to be satisfied at once. A female can only conceive if she is not already pregnant, is older than the minimum reproductive age, has satiety above a minimum threshold, is off her reproductive cooldown, and perceives a male nearby. The satiety requirement is what couples reproduction to food: in a poor environment, fewer females clear the bar and the birth rate falls on its own, without us having to script it. Once pregnant, the female carries for a gestation period of about one year, and when it elapses she gives birth and enters a multi-year cooldown before she can reproduce again, which reproduces the long inter-birth intervals that are characteristic of the species.

The litter is not a fixed number. Its size is drawn from a normal distribution around a calibrated mean (close to two cubs), clamped to be at least one. Each cub is then subjected to an infant-mortality probability at birth, representing perinatal losses, so not every cub in the litter is actually added to the live population. This two-step birth process, variable litter size followed by per-cub survival, gives a more realistic recruitment signal than simply adding a fixed number of healthy cubs per mother.

\subsection{Avoiding artificial birth waves at founding}

There is a subtle but important detail in how the founding population is created. If every founding female of reproductive age started the simulation in exactly the same reproductive state, they would all come off cooldown or finish gestation at the same time, producing an artificial wave of synchronised births that has nothing to do with ecology and everything to do with how we set up the run. To avoid this, the founding reproductive females are staggered: a fraction of them start already pregnant with a random remaining gestation, and the rest start with a random remaining cooldown. The founding ages themselves are not drawn uniformly either. A flat age spread from zero to twenty-five years would have a mean around twelve and a half, which is roughly double that of a real brown bear population and over-represents middle-aged bears, again triggering a boom-then-crash transient. Instead, founders are sampled from a young-skewed, truncated-exponential age pyramid with a mean of about seven years, which is much closer to a real stable age distribution \cite{Swenson1997,GrimmRailsback2005}. These two choices together make the population start near its stationary structure rather than lurching into it.

\section{Environment representation}

The world is a square grid, and each cell carries a terrain type together with a food value and a danger value, both living in the interval $[0, 1]$, plus explicit occupancy information. The terrain types are forest, field, village, road, mountain and a blocked boundary type. When a random map is generated, the terrain mix is controlled by configurable percentages: in the current settings about a quarter of the map is forest, more than half is field, a tenth is mountain, and small fractions are road and village. Each terrain type then has its own average food and average danger, with a bounded random variance on top so that no two cells of the same type are identical. Forest is the rich, safe habitat (high average food, low danger), villages are food-rich but extremely dangerous, roads are poor and dangerous, and fields and mountains sit in between. This is a compact way of encoding the central ecological tension of the whole problem: the safest food is in the forest, but the easiest, most concentrated food is exactly where it is most dangerous to be, near people.

Food is not static. Bears reduce the food on a cell when they eat there, and the environment regrows food over time, restoring a random number of cells each tick. The amount eaten per tick and the amount regrown per tick are both parameters, and the ratio between them effectively sets the carrying capacity of the landscape. This is precisely the lever that the anthropogenic-food scenario pulls: increasing the regrowth rate over the years is how we represent more food becoming available to bears, whether through supplemental feeding, crops, or waste near settlements.

The environment also exposes a percept for each agent. The percept contains the bear's current cell with its food, danger and local bear count, the eight neighbouring cells with the same information plus a blocked flag, and a single boolean saying whether a male is nearby, which the females use for the reproduction check. The perception is strictly local; an agent never sees beyond its immediate ring of cells, which keeps the model faithful to the idea that bears act on local information.

Finally, the implementation can either generate a random map from the terrain percentages or load a pre-processed map grid from a file. This second capability is what allows the project to move beyond purely synthetic landscapes toward geographically structured experiments, where the grid reflects a real habitat layout rather than a statistical mix.

\section{Module-level architecture and the action--effects pattern}

To keep the implementation understandable and to let the ecological logic, the execution policy and the experiment workflow evolve independently, the code is separated into modules with clear responsibilities. The \emph{BearAgent} module holds the internal state and the decision logic. The \emph{BearEnvironment} module builds percepts and applies state transitions. The \emph{BearSimulation} driver owns the tick loop, the phase boundaries and the scheduling. A dedicated effects container decouples deciding from mutating. And an \emph{Experiments} module handles run configuration, metrics output, schedules and calibration tools. This is the same separation-of-concerns principle that runs through the whole thesis: keep things that change for different reasons in different places.

The action--effects pattern is worth describing on its own because it is what makes the parallel-plan / serial-commit idea actually work. When an agent decides on an action, the action is not applied directly to the world. Instead, each action object, when asked to contribute to the current step, writes its \emph{intent} into a shared per-tick \emph{BearActionEffects} container. A move writes a move intent keyed by agent id; a death writes a die intent together with its cause; a birth writes a birth intent naming the mother; eating writes a food-consumption intent for a cell; and a counter is kept for human-conflict exposure. Nothing in the world has changed yet at this point, we have only gathered a description of what everybody wants to do.

The commit phase then resolves these intents in a fixed, deliberate order on the driver thread: first deaths are applied and the dead bears are removed, then births add new cubs after litter sampling and infant mortality, then food reductions are applied to the eaten cells, then move intents are resolved, then conflict exposure is recorded, and finally food is partially regrown. Applying intents in a fixed order is what gives the model deterministic semantics: there is never a question of which of two simultaneous changes ``won'', because the order is always the same. A pleasant side effect of centralising every change through this container is observability, because all the intents pass through one place, it is trivial to count births, deaths by cause, conflict events and so on for each tick and export them, which is the backbone of the entire results chapter.

One detail about movement is worth being honest about: the model does not enforce that a cell can hold only one bear. Move intents are grouped by their target cell and every claimant is moved there. Bears can therefore share a cell, and competition between them is expressed softly, through the crowding penalty in the movement score and through the shared food on the cell, rather than through a hard collision rule. This is a simplification, but it is a reasonable one at the scale of the question we are asking.

\section{Execution architecture and concurrency}

The simulation uses parallel planning and serial commit as its core execution strategy, and the implementation leans on a modern Java feature to do it cheaply. During planning, one small task is created per agent. Each task fetches that agent's percept and asks the agent to decide its action, and all of these tasks are submitted to an executor that uses \emph{virtual threads} (\texttt{newVirtualThreadPerTaskExecutor}). Virtual threads are extremely lightweight, so we can afford to create one task per bear even when there are several thousand bears, without worrying about exhausting operating-system threads. The driver then collects the results into an ordered map keyed by agent id and hands them to the commit phase.

The reason this works without any locks on the world is precisely the plan/commit separation discussed earlier. During planning, the workers only \emph{read} the world to build percepts and they only \emph{write} to the agent's own private state, never to shared structures. The shared world is touched only in the commit phase, which runs on a single thread. From a software-engineering point of view, this design gives three concrete advantages: safe concurrency for the expensive decision-making, a consistent and reproducible world transition per tick, and a single natural place to plug in metrics, schedule hooks and reproducibility controls.

There is a trade-off here, of course: the serial commit is a hard synchronisation barrier. Every tick, all the parallel work has to finish and then funnel through one thread, and the loop cannot advance until the commit is done. At very large scales the commit can become the bottleneck, and, as we will argue in the computational discussion, this same barrier is exactly what makes distributing a \emph{single} run across machines unattractive. In the current version this trade-off is accepted on purpose, because correctness and interpretability are prioritised over squeezing out the last bit of low-level performance.

\section{Reproducibility pipeline}

Probably the strongest part of the whole project is the reproducibility infrastructure, due to the fact that it is what turns the simulator from a nice demonstration into something that can actually back up a scientific claim. The principle behind it is simple: every reported number must be traceable back to a scenario definition and a random seed.

The trickiest part of reproducibility in a parallel simulation is randomness. If all agents drew from one shared random generator, then the order in which the threads happened to run would change the sequence of draws, and the same seed would give different results on different machines or even different runs. The \emph{RngSupport} module solves this with per-agent random streams. In seeded mode, each agent gets its own \emph{Random} instance whose seed is derived from the master seed and the agent id through a SplitMix64-style mixing function. Because the stream depends only on the pair (master seed, agent id), a given agent always sees exactly the same sequence of random numbers regardless of which worker thread happens to execute its decision. A separate environment-level generator, used only from serial code such as map generation and the commit phase, handles all the world-level randomness. The per-agent generators are kept in a concurrent map so that threads planning different agents never contend when a new agent's stream is created. When no seed is set, everything falls back to \texttt{ThreadLocalRandom}, so ordinary interactive runs pay no penalty for machinery they do not need.

Around this sits the rest of the pipeline. A run configuration object captures the scenario id, the run id, the seed, the replicate index and the maximum number of ticks. A per-tick metrics recorder samples the population, births, deaths by cause, satiety, age structure and conflict events, and the founding cohort is recorded before any tick runs so that the initial gender split and founding-pregnancy count are visible too. The \emph{BatchRunner} ties it together: it runs several seeded replicates of each scenario, applies any parameter overrides to the global settings before a run and restores them afterwards, writes one per-tick CSV per replicate, and writes a metadata CSV and a summary CSV with the fit metrics. Every result in the next chapter comes out of this pipeline.

\section{Schedule-driven scenario modelling}

A static model cannot represent the question we are asking, because the phenomenon itself is time-dependent: protection policies came into force in particular years, hunting rules changed, food availability shifted over the period. The model therefore supports time-indexed parameter changes through a \emph{ParameterSchedule} and a \emph{ScheduleApplier}. A schedule is simply a list of entries, each saying ``at this simulation year, set this parameter to this value''. As the simulation clock crosses each scheduled year, the applier mutates the corresponding field.

The implementation does this through reflection on the static \emph{Settings} class, which is an elegant point I am quite happy with: a schedule can target \emph{any} settings field by name, without the simulation code having to know in advance which parameters a scenario will want to change. The applier looks up the named field, checks that it is static, parses the new value to the field's type, and sets it. Crucially, the first time it touches a field it snapshots the original value, and it records every change it makes as an applied event with the tick, the year, the parameter, the old value and the new value. At the end of the run it restores every field it ever touched, so one run can never silently contaminate the next, and the recorded events become a per-run audit trail that is written out alongside the results. This is how scenarios such as ``danger mortality drops from year two'' (the 2007 EU protection driver), ``food regrowth increases stepwise in years five, ten and fifteen'' (anthropogenic food), or a layered combination of several such changes are expressed, all without writing any scenario-specific code in the simulation loop itself.

\section{Skipping the initialization transient}

A practical problem with starting every run from freshly generated random founders is the initialization transient: even with the young-skewed age pyramid and the staggered reproductive states described earlier, the artificial founding population needs a few simulated years to settle into a realistic spatial and demographic structure, and during that settling the trajectory is not yet meaningful. Running every scenario through that warm-up over and over is wasteful and, worse, it injects a wobble into the early years that we then have to explain away.

The project addresses this in two complementary ways. The first is snapshotting: the simulation can be run once until the population has spun up, then a \emph{BearSnapshot} of the whole world and population is written to disk, and later runs can restore that snapshot instead of generating founders, so they start from an already-settled state with the clock reset to zero. The second is the \emph{BalancedInitializer}, which goes a step further: it reads a characterization of the measured stable state (age distribution, satiety, sex, reproductive state, per-habitat placement) and generates a fresh founding population directly from those distributions, while also initialising the map's food to its drawn-down level rather than a full larder. Both approaches remove the transient; the difference is that the snapshot restores one specific frozen population exactly, while the balanced initializer samples a new but statistically equivalent one. For the campaign in the next chapter this matters a great deal, because it means the scenarios all start from the same realistic baseline and the differences we measure between them are due to the drivers, not to leftover warm-up noise.

\section{Performance instrumentation}

So that the computational discussion can be based on evidence rather than impression, the simulation contains a detailed timing instrumentation, active when benchmarking is enabled. The loop time is split into the planning time and the commit time, and each of those is decomposed further. On the planning side we record the time spent building the worker tasks, invoking them all, and collecting their futures, plus the aggregate worker CPU time spent on percept generation and on decision-making (measured per worker and therefore non-additive, since the workers run concurrently). On the commit side we record, separately, the time spent contributing actions into the effects container, applying deaths, applying births, applying food reductions, resolving moves, regrowing food and recording performed actions. There is also an explicit accounting of ``unaccounted'' overhead in each phase, so the numbers add up honestly.

This level of detail is what lets us say something precise instead of vague in the next chapter. Rather than claiming that one execution mode is ``faster'', we can point to exactly where the time goes, and, in particular, show how much of each tick is spent in the serial commit barrier, which is the quantity that decides whether distributing a single run could ever pay off.

\section{Design limitations and current constraints}

It is worth being honest about the current limits of the architecture, and these are best understood as the boundary of where the approach currently applies rather than as defects.

First, many parameters live in a single static \emph{Settings} class. This is convenient and it is what makes the reflection-based scheduler so simple, but global mutable state has to be handled carefully: the schedule applier's snapshot-and-restore discipline exists precisely to stop one run leaking into another, and strict isolation between runs depends on that discipline being respected. Second, the tick barrier imposes strict phase boundaries; this is what guarantees determinism, but it is also what limits how far a single run can be distributed across machines. Third, several ecological processes are intentionally abstracted, food is a scalar field rather than distinct food types, conflict is a proxy count of bears on village and road cells rather than a modelled human response, and there is no genetic or individual-quality variation between bears beyond age and sex. These are conscious simplifications appropriate to the scope of this thesis and to the comparative question being asked, and each of them is a natural direction for future extension.

\section{Multithreading, and why distribution was not applied}

It is worth being completely clear about this, because it would be easy to overstate it: the model is executed in a multithreaded way on a single machine, and a distributed execution was \emph{not} implemented anywhere in the project. This is a deliberate decision rather than an omission, and the reasoning follows directly from everything above. I want to explain why distribution simply did not make sense here, instead of bolting on a distributed mode just so the project could carry the label.

Each agent's per-tick work, that is building a local percept, scoring eight neighbouring cells and picking an action, is very small, on the order of microseconds. At the same time the model is tightly tick-synchronised: every tick ends in a serial commit barrier that all agents must reach before anyone can advance. If one tried to spread the agents of a single run across several machines, every single tick would require shipping each machine's intents back to a coordinator, resolving them, and shipping the updated world back out, all before the next tick could begin. The communication and synchronisation cost of doing that, tens of thousands of times over a sixteen-year run, would dwarf the tiny amount of useful work it parallelises, so a distributed single run would almost certainly be \emph{slower} than the multithreaded one, not faster. This is a well-known pattern in the parallel and distributed simulation literature: fine-grained, frequently synchronised workloads do not distribute well, whereas coarse-grained ones do \cite{Fujimoto2000}.

The right way to think about parallelism here is to match it to the granularity that actually exists in the problem:

\begin{itemize}
    \item multithreading, with the lightweight virtual-thread planning stage, is what the project uses to accelerate the agent decisions \emph{inside} one run on a single machine, where the threads share memory and the commit barrier costs essentially nothing;
    \item distribution would only ever belong at the level of \emph{many independent runs}, different scenarios and different replicate seeds, since those runs share nothing and never synchronise with each other.
\end{itemize}

That batch level is the embarrassingly parallel part of the problem, and it is the only place distribution could have actually paid off. But in practice it was not needed: the whole experiment campaign is a manageable number of independent seeded runs, and a single multi-core machine ran them in acceptable time, so there was no reason to add the complexity and the failure modes of a distributed setup. The honest summary, then, is that the project is a multithreaded agent-based model whose design was \emph{analysed} for distribution and found not to warrant it, which is itself one of the more useful technical findings of the thesis, because it pushes back on the common assumption that more machines automatically means a faster simulation.

\section{Chapter summary}

This chapter described the proposed approach in full, from the modelling assumptions and the agent's decision pipeline, through the environment and the action--effects commit logic, to the concurrency model, the reproducibility pipeline, the schedule-driven scenarios, the spin-up machinery and the performance instrumentation. The recurring theme has been the plan/commit separation and the decision to keep ecological logic, execution policy and experiment workflow apart from one another. The next chapter puts all of this to work: it defines the evaluation protocol and reports both the ecological results of the 2005--2021 campaign and the computational findings that justify the multithreaded single-machine execution and the reasoned decision not to distribute.

\chapter{Evaluation and validation}
\label{ch:evaluation}

\section{Evaluation strategy}

The evaluation in this thesis is centered on ecological explanatory fit for the 2005--2021 window, using scenario-based counterfactual analysis.

The reference annual trajectory is the Romanian national series (2005--2021), from 5800 bears in 2005 to 7500 bears in 2021. This implies an observed cumulative increase of 29.31\% and an average annual growth of approximately 1.62\%/year.

The simulations are compared against this trajectory with:

\begin{itemize}
    \item MAE (mean absolute error),
    \item RMSE (root mean square error),
    \item end-of-window gap to 2021 reference,
    \item cumulative growth and annualized growth inside the same window.
\end{itemize}

\section{Experimental protocol and scenarios}

All scenarios were run with three replicates and then post-processed into annual aggregates (year 0 to year 16, mapped to 2005 to 2021). The campaign includes:

\begin{itemize}
    \item baseline,
    \item EU protection driver (2007),
    \item hunting-ban driver (2016),
    \item intervention-cull driver (2021),
    \item habitat-loss driver,
    \item anthropogenic-food driver,
    \item combined realistic schedule.
\end{itemize}

\begin{table}[htbp]
\centering
\caption{Scenario definitions used in the final campaign.}
\label{tab:ch4_scenario_matrix}
\begin{tabular}{|l|p{8.5cm}|}
\hline
Scenario & Main schedule effect \\
\hline
Baseline & No time-varying policy/habitat driver beyond baseline settings. \\
EU protection 2007 & Danger mortality reduced from simulation year 2. \\
Hunting ban 2016 & Danger mortality reduced from simulation year 11. \\
Intervention cull 2021 & Danger mortality increased from simulation year 16. \\
Habitat loss & Forest food reduced stepwise (years 3, 9, 14). \\
Anthropogenic food & Food regrowth increased stepwise (years 5, 10, 15). \\
Combined realistic & Layered schedule: protection + food increase + 2016 ban + 2021 cull. \\
\hline
\end{tabular}
\end{table}

\section{Ecological results}

\subsection{Quantitative summary}

Table~\ref{tab:ch4_ecology_summary} shows the core quantitative comparison against the 2005--2021 reference trajectory.

\begin{table}[htbp]
\centering
\caption{Ecological validation metrics (annual aggregates, 2005--2021).}
\label{tab:ch4_ecology_summary}
\begin{tabular}{|l|c|c|c|c|c|}
\hline
Scenario & End pop. & Growth (\%) & CAGR (\%/y) & MAE & RMSE \\
\hline
Baseline & 6020.67 & 2.57 & 0.16 & 360.32 & 539.66 \\
EU protection 2007 & 6103.33 & 3.98 & 0.24 & 342.60 & 511.44 \\
Hunting ban 2016 & 6068.67 & 3.39 & 0.21 & 348.76 & 522.58 \\
Intervention cull 2021 (16y) & 6020.67 & 2.57 & 0.16 & 360.32 & 539.66 \\
Habitat loss & 6020.67 & 2.57 & 0.16 & 360.32 & 539.66 \\
Anthropogenic food & 7175.00 & 22.24 & 1.26 & 261.25 & 336.27 \\
Combined realistic & 7194.67 & 22.58 & 1.28 & 278.64 & 344.65 \\
\hline
Reference (observed) & 7500.00 & 29.31 & 1.62 & 0.00 & 0.00 \\
\hline
\end{tabular}
\end{table}

Main observations:

\begin{enumerate}
    \item Baseline and policy-only schedules underpredict growth strongly in this calibration.
    \item The anthropogenic-food schedule yields the best fit in this campaign (lowest MAE and RMSE).
    \item The combined realistic schedule is very close to anthropogenic food, but slightly worse on MAE/RMSE.
    \item The 2021 intervention-cull schedule has negligible effect within the 2005--2021 comparison window because it is activated at the end of the window.
\end{enumerate}

Relative to baseline, MAE drops by about 27.5\% under the anthropogenic-food scenario and by about 22.7\% under the combined scenario.

\subsection{End-of-window interpretation}

At year 16 (mapped to 2021), the reference value is 7500 bears.

\begin{itemize}
    \item Baseline gap: 1479.33 bears below reference.
    \item Anthropogenic-food gap: 325.00 bears below reference.
    \item Combined realistic gap: 305.33 bears below reference.
\end{itemize}

This indicates that, in the current parameterization, increased food availability is the dominant mechanism needed to approach the observed 2021 level.

\subsection{Mortality-structure summary}

Across scenarios, mean annual total deaths decrease from approximately 512.76 (baseline) to about 491.47--491.61 in the two best-fit scenarios. The danger-death share also decreases in the combined schedule relative to baseline, consistent with lower danger mortality settings in part of the timeline.

\section{Figure placeholders for final insertion}

The following placeholders are intentionally left as boxes so final plots can be inserted after figure export.

\begin{figure}[htbp]
    \centering
    % Suggested final file: figures/ch4_population_all_scenarios.pdf
    \fbox{\rule{0pt}{2.2in}\rule{0.85\textwidth}{0pt}}
    \caption{Observed reference trajectory (2005--2021) versus all simulated scenario trajectories.}
    \label{fig:ch4_population_all}
\end{figure}

\begin{figure}[htbp]
    \centering
    % Suggested final file: figures/ch4_population_top3.pdf
    \fbox{\rule{0pt}{2.2in}\rule{0.85\textwidth}{0pt}}
    \caption{Focused comparison of baseline, anthropogenic-food, and combined-realistic scenarios against observed data.}
    \label{fig:ch4_population_top3}
\end{figure}

\begin{figure}[htbp]
    \centering
    % Suggested final file: figures/ch4_residuals_by_year.pdf
    \fbox{\rule{0pt}{2.2in}\rule{0.85\textwidth}{0pt}}
    \caption{Year-by-year residuals (simulated minus observed) for selected scenarios.}
    \label{fig:ch4_residuals}
\end{figure}

\begin{figure}[htbp]
    \centering
    % Suggested final file: figures/ch4_deaths_components_selected.pdf
    \fbox{\rule{0pt}{2.2in}\rule{0.85\textwidth}{0pt}}
    \caption{Deaths by cause (danger, starvation, old age) for baseline and best-fit scenarios.}
    \label{fig:ch4_death_components}
\end{figure}

\begin{figure}[htbp]
    \centering
    % Suggested final file: figures/ch4_births_vs_deaths_selected.pdf
    \fbox{\rule{0pt}{2.2in}\rule{0.85\textwidth}{0pt}}
    \caption{Births versus total deaths over time for baseline and best-fit scenarios.}
    \label{fig:ch4_births_deaths}
\end{figure}

\section{Computational validation note}

The current chapter uses the newly generated ecological campaign data from the Results directory. Detailed runtime benchmarking (sequential vs multithreaded vs distributed) is discussed in Chapter 3 and can be appended here with the same table/figure pattern if a dedicated benchmark rerun is performed on the final hardware.

\section{Interpretation and validity limits}

The most important conclusion from this campaign is comparative: in the tested schedules, food-access increase explains substantially more of the 2005--2021 rise than policy-only danger-mortality shifts.

Two limits are important:

\begin{enumerate}
    \item The folders marked as weird results indicate the need for an additional reproducibility check before publication-grade claims about driver ranking.
    \item The intervention-cull effect should be evaluated on a longer post-2021 projection horizon, because activation at year 16 leaves almost no in-window response time.
\end{enumerate}

Even with these limits, the campaign is strong enough to support a final thesis-level conclusion about which mechanisms are necessary to approach observed growth in the validated window.

\chapter{Discussion, conclusion and future work}
\label{ch:discussion_conclusion}

\section{Discussion}

Taken together, the results of the final campaign give a fairly clean ranking of explanatory power inside the 2005--2021 validation window, and it is a ranking that I find convincing precisely because the different pieces of evidence all agree with one another instead of pulling in different directions.

The baseline and the policy-only schedules, the 2007 EU protection and the 2016 hunting-ban reductions in danger mortality, improve the fit only marginally relative to the baseline. The schedules that increase food access over time, by contrast, produce a clearly stronger match to the observed trajectory. The single best scenario in the run set is anthropogenic food, with an MAE of 261.25 against 360.32 in the baseline, an RMSE of 336.27 against 539.66, and an end population of 7175 against the 7500 reference for 2021. The combined realistic schedule comes very close (MAE 278.64, RMSE 344.65, end population 7194.67) but does not actually beat anthropogenic food on aggregate error. That last point is more interesting than it first looks: the combined schedule contains the food increase \emph{and} the policy drivers, so the fact that adding the policies on top of the food does not improve the fit tells us that, in this parameterization, the food pathway is doing essentially all of the explanatory work, while the additional policy components contribute little or partly offset one another.

The reason behind this ranking becomes obvious once we look at the mortality structure. In the model, the overwhelming majority of deaths are from starvation, on the order of four hundred a year against only a few dozen from danger. The policy drivers act solely on danger mortality, which is the small component, so even a large relative cut to it can only recover a handful of bears. The food driver acts on the dominant component directly: more food means fewer bears starve, and at the same time more females clear the satiety requirement and reproduce. It pushes on the biggest lever in the system from both the death side and the birth side at once, which is exactly why it is the only mechanism in the campaign able to close most of the gap to the observed level. The year-by-year trajectory reinforces this with a believable lag: the food increase begins in year~5 but the curves do not visibly separate until around year~8, because more food converts into more bears only through the slow chain of better body condition, then more births, then surviving cubs, a chain that carries the long reproductive cycle of the species inside it.

A second point worth dwelling on is horizon sensitivity, because it is easy to misread. The 2021 intervention-cull schedule is activated at simulation year~16, the very end of the comparison window, so within 2005--2021 it has essentially no time to act and comes out nearly identical to the baseline. This should not be read as ``culling does nothing''. It is purely a consequence of timing: the driver fires too late to express itself inside the validated window, and measuring its real effect would require extending the simulation past 2021. This is a good illustration of why a schedule-driven agent-based model is valuable for this kind of policy question in the first place: it lets us reason not only about whether a driver causes growth or decline, but about \emph{when} a driver is applied and how long it needs to take effect.

\section{Conclusion}

This thesis set out to answer which mechanisms can best explain the increase of the Romanian brown bear population over 2005--2021, and to do so with a model that is reproducible and whose computational behaviour is understood rather than assumed. It delivers both a methodological contribution and a data-backed empirical conclusion.

\subsection*{Methodological contribution}

\begin{enumerate}
    \item A spatial, agent-based brown bear simulator with explicit demographic and environmental mechanisms, age- and sex-structured mortality, satiety-gated reproduction with gestation and cooldown, variable litters with infant mortality, seasonal hibernation, subadult-male dispersal, and a food/danger landscape that the bears both consume and respond to.
    \item A reproducible scenario pipeline built directly into the code: per-agent seeded randomness that survives parallel execution, time-indexed parameter schedules applied and audited through reflection, seeded replicates, spin-up via snapshots and balanced initialization, annual post-processing, and quantitative fit metrics.
    \item A practical and somewhat unusual computational result for this class of model: the model is multithreaded on a single machine, and a distributed execution was deliberately not implemented, because the timing analysis shows that a fine-grained, tick-synchronised run would lose more to communication and synchronisation overhead than it could gain. Distribution would only make sense around a batch of independent runs, and at this scale a single machine made even that unnecessary.
\end{enumerate}

\subsection*{Empirical conclusion for 2005--2021}

Within the tested scenarios and the current calibration, population growth close to the observed Romanian trajectory requires an increase in food availability over time. Policy-only adjustments to danger mortality are not, by themselves, sufficient to reproduce the magnitude of the rise. In short, the campaign supports a multi-driver interpretation in which food-access dynamics are both necessary and dominant in this version of the model. This is a comparative statement, made carefully and within stated assumptions, and it is exactly the kind of statement the scope of the thesis was designed to be able to support.

\section{Limitations}

I would rather state the limitations plainly than have them found, because each of them marks the boundary of where the present work applies, not a flaw that undermines it.

\begin{enumerate}
    \item Two of the result folders, the anthropogenic-food and the combined-realistic outputs, are flagged in the data directory as needing a reproducibility recheck before they are used for publication-grade attribution, because of a possible pipeline contamination between two closely related scenarios. The comparative ranking is robust, but the precise error figures for those two scenarios should be re-confirmed under a strict, isolated protocol.
    \item The validation uses a single national annual reference trajectory. Spatially explicit, local validation, checking that the model gets the right bears in roughly the right places, is outside the current scope, and a good fit to the national curve does not by itself guarantee a correct spatial pattern.
    \item Some drivers, especially the late-window 2021 cull, are activated too near the end of the window to express their full effect, so their real impact can only be quantified with an extended horizon.
    \item Several ecological processes are intentionally abstracted: food is a single scalar field rather than distinct food types with their own seasons, human--bear conflict is a proxy count of bears on village and road cells rather than a modelled human response, and there is no individual-quality or genetic variation between bears beyond age and sex.
\end{enumerate}

These limitations constrain the strength of the claims, but they do not invalidate the comparative ranking that the campaign produced, and most of them point directly at the next steps.

\section{Future work}

There are five directions in which I would take this work further, and they follow naturally from the limitations above.

\begin{enumerate}
    \item Run a post-2021 extension campaign, for example to 2026 or 2030, specifically to measure the delayed effect of the intervention-cull driver and of any other late-window policy, which the current window cannot capture.
    \item Recheck and rerun the anthropogenic-food and combined scenarios under a strict, fully isolated reproducibility protocol, to rule out the possible pipeline contamination and to confirm the precise error figures.
    \item Add uncertainty envelopes, confidence bands across replicates, directly onto the final trajectory figures, so that the comparison between scenarios is shown with its variability rather than as bare lines.
    \item Extend the validation with additional anchors, for example independent genetic population estimates, while keeping the annual-series evaluation separate, so that the model is checked against more than one kind of data.
    \item Add a dedicated benchmark appendix tied to one final, frozen code-and-data snapshot, so that the ecological and computational conclusions are both anchored to the exact same version of the system, and so that the multithreaded scaling (wall-clock time against agent count and core count) can be reported as hard numbers on specific hardware. If the experiment scale ever grew enough to justify it, the same appendix would be the place to prototype and measure a distributed batch layer and confirm in practice the cost-benefit argument made here.
\end{enumerate}

Beyond these, the natural longer-term extensions are to enrich the environment, distinct food types with their own seasonality, a real loaded habitat grid for Romania's forest range instead of a statistical terrain mix, and to model the human side of conflict explicitly rather than as a proxy count.

\section{Final remarks}

The thesis is in a complete state: the numerical evidence is integrated, the scenario ranking is explicit, and the result figures are in place. What remains before submission is mechanical rather than conceptual, that is cross-checking every table value against the final exported CSVs and a last pass over the narrative for consistency of captions, references and the links between the schedules and their results. The substance of the argument, ecological and computational alike, is already on the page.

%\addcontentsline{toc}{chapter}{Concluzii}
%\addcontentsline{toc}{chapter}{Conclusions}

\bibliography{references}

\end{document}
